% Generated draft: requires mathematical and source review.
\documentclass[11pt]{article}
\usepackage[margin=25mm]{geometry}
\usepackage{fontspec}
\setmainfont{Latin Modern Roman}
\usepackage{amsmath,amssymb,mathtools,unicode-math}
\setmathfont{Latin Modern Math}
\usepackage{enumitem,needspace,xurl,hyperref}
\setlength{\parindent}{0pt}
\setlength{\parskip}{5pt}
\setlength{\emergencystretch}{3em}
\newcommand{\N}{\mathbb{N}}
\newcommand{\Z}{\mathbb{Z}}
\newcommand{\Q}{\mathbb{Q}}
\newcommand{\R}{\mathbb{R}}
\newcommand{\C}{\mathbb{C}}
\newcommand{\F}{\mathbb{F}}
\newcommand{\A}{\mathbb{A}}
\newcommand{\PP}{\mathbb{P}}
\newcommand{\E}{\mathbb{E}}
\newcommand{\eps}{\varepsilon}
\newcommand{\dd}{\,\mathrm{d}}
\begin{document}
\section*{UnsolvedMath ID: 30002514}
\textbf{Source problem code:} catalogue entry
\subsection{Definitions and mathematical statement}
The supplied problem asks for a description of the orbits of \(G=\operatorname{Gal}(\overline{\Q}/\Q)\) on isotopy classes of finite planar trees. Its stated equivalence with Shabat polynomials requires conventions about vertex colours, polynomial equivalence, and allowed critical values that the record does not specify completely.

A precise standard dessin formulation can be given as follows. A bicoloured plane tree is a finite connected acyclic graph with at least one edge, a cyclic ordering of the edges incident to each vertex, and a colouring of its vertices by \(0\) and \(1\), with adjacent vertices having different colours. Trees are considered up to graph isomorphisms preserving the cyclic orderings and the colours. The cyclic orderings encode an embedding in an oriented plane; no root or vertex labels are imposed.

For each integer \(n\geq 1\), set
\[
\mathcal S_n=\left\{P\in\overline{\Q}[z]:\deg P=n,\quad P'(c)=0\Longrightarrow P(c)\in\{0,1\}\text{ for every }c\in\overline{\Q}\right\}.
\]
Thus the set of finite critical values is required to be contained in \(\{0,1\}\). Define \(P\sim Q\) if
\[
Q(z)=P(az+b)\qquad\text{for some }a\in\overline{\Q}^{\times},\ b\in\overline{\Q}.
\]
Under the usual dessin correspondence, the classes in \(\mathcal S_n/\!\sim\) correspond to bicoloured plane trees with \(n\) edges: the embedded tree is \(P^{-1}([0,1])\subset\C\), with vertices over \(0\) and \(1\) carrying their respective colours.

For \(P(z)=\sum_{j=0}^n a_jz^j\) and \(\sigma\in G\), define
\[
P^\sigma(z)=\sum_{j=0}^n\sigma(a_j)z^j,
\qquad \sigma\cdot[P]=[P^\sigma].
\]
This defines an action on each \(\mathcal S_n/\!\sim\), and hence, through the correspondence, on bicoloured plane trees. With these conventions, the problem is to describe completely the orbit partition
\[
\mathcal S_n/\!\sim\;=\coprod_{\mathcal O\in G\backslash(\mathcal S_n/\!\sim)}\mathcal O
\qquad(n\geq1).
\]
Equivalently, one seeks a complete characterization of when two associated trees belong to the same orbit; in polynomial terms the relation is
\[
[P]\text{ and }[Q]\text{ lie in the same orbit}
\quad\Longleftrightarrow\quad
\exists\sigma\in G\ \exists a\in\overline{\Q}^{\times}\ \exists b\in\overline{\Q}\quad
Q(z)=P^\sigma(az+b).
\]

\emph{Unresolved source conventions.} If the intended trees have no chosen bipartition colouring, the corresponding polynomial classes must additionally identify \(P\) with \(1-P\). This involution commutes with the Galois action, so an action also descends to that quotient. The supplied record does not distinguish these two classification problems. Also, ``whose only critical values are \(0\) and \(1\)'' may require both values to occur, in which case \(\mathcal S_n\) must instead require
\[
\{P(c):P'(c)=0\}=\{0,1\}.
\]
The containment convention includes degree-one polynomials and polynomials with just one finite critical value; the equality convention excludes them. A one-vertex tree with no edges is not covered by the nonconstant-polynomial formulation, and its intended inclusion is unspecified. Thus the formulation above states an explicit standard interpretation while leaving these source qualifications unresolved.
\subsection{Short English statement}
Describe all absolute Galois orbits of finite plane trees through their associated Shabat polynomials. A precise version uses bicoloured plane trees and polynomials modulo affine changes of the input variable, with Galois acting on coefficients. The source leaves unresolved whether colours are forgotten, whether both critical values must occur, and whether the tree with no edges is included.
\subsection{Sources}
\begin{itemize}
\item Supplied local UnsolvedMath registry record, permanent ID \(30002514\), title \emph{Galois Orbits of Finite Planar Trees}, exact source problem code \texttt{catalogue entry}, in the Oberwolfach Reports Open Problems collection. The record gives proposed year \(2014\) and records no named proposer.
\item Supplied report citation: \emph{New Trends in Teichm\"uller Theory and Mapping Class Groups} (2014), Oberwolfach Reports, pp.~395--452, DOI \texttt{10.4171/owr/2014/07}. Supplied literature-source label: Oberwolfach Report 07/2014. \url{https://doi.org/10.4171/owr/2014/07}. The record describes these reports as documenting mathematical workshops and their open-problem discussions; no individual problem page or speaker attribution is supplied.
\item Supplied literature-source label: Schneps--Shabat, plane-tree orbits. \url{https://www.cambridge.org/core/books/abs/grothendieck-theory-of-dessins-denfants/on-the-classification-of-plane-trees-by-their-galois-orbit/FD4DB49548FCE9389A8D980C90A757BF}.
\item Supplied literature-source label: Dessins database. \url{https://irma.math.unistra.fr/~guillot/research/dessins_database/index.html}.
\item Literature assessment dated exactly \texttt{2026-08-21}: the supplied curation labels the problem \texttt{partially\_solved} and states that Galois invariants and substantial classifications of plane-tree dessins are known, that databases compute many orbits, and that no complete description of the Galois orbits of finite planar trees is available. This is supplied dated curation derived from the OWR record, not an independently verified assessment. The supplied background expressly requires independent verification of time-sensitive status and cited claims.
\item The registry labels the statement \texttt{corrected\_verified} and says that notation, typography, and wording were normalized without changing the mathematical content. These are supplied registry assertions, not independent verification of this draft or of the cited sources.
\end{itemize}
\end{document}
